% SPDX-License-Identifier: CC-BY-SA-4.0
% Three slide patterns: a statement, a six-beat chase, and a working proof.
% Input from atlas.tex, after \begin{document}. TikZ and booktabs are loaded.
% Every diagram has the same coordinates and explicit bounding box. A fixed
% overlay area also keeps replacement prose from moving the picture.

% Arguments: cop vertex, robber vertex, remaining vertices, solid edges.
% The remaining vertices are the component of T-c containing r. At capture
% we show the empty set by convention. Edges behind the cop stay dashed;
% lettered circle/square pieces also distinguish players in the plain deck.
\newcommand{\AtlasPursuitTree}[4]{%
  \begin{tikzpicture}[x=1.15cm,y=1cm,
    edge/.style={line width=.8pt,draw=AtlasInk},
    faded/.style={edge,draw=AtlasInk!35!AtlasPaper,dashed},
    cop/.style={circle,draw=AtlasRubric,fill=AtlasRubric,text=AtlasPaper,
      minimum size=6mm,inner sep=0pt,font=\small},
    robber/.style={rectangle,draw=AtlasInk,fill=AtlasPaper,text=AtlasInk,
      line width=1pt,minimum size=6mm,inner sep=0pt,font=\small}]
    \path[use as bounding box] (-.45,-1.9) rectangle (4.95,1.85);
    \begin{scope}[overlay]
      \coordinate (a) at (0,0);
      \coordinate (b) at (1.1,0);
      \coordinate (c) at (2.2,0);
      \coordinate (d) at (3.3,.9);
      \coordinate (e) at (4.5,.9);
      \coordinate (f) at (1.1,1.3);
      \coordinate (g) at (1.1,-1.2);
      \coordinate (h) at (2.2,-1.2);
      \foreach \u/\v in {a/b,b/c,c/d,d/e,b/f,b/g,c/h}
        \draw[faded] (\u)--(\v);
      \foreach \u/\v in {#4}
        \draw[edge] (\u)--(\v);
      \foreach \v in {a,b,c,d,e,f,g,h} {
        \fill[AtlasInk!35!AtlasPaper] (\v) circle (1.8pt);
      }
      \foreach \v in {a,d,e,g,h}
        \node[below=3.5mm,font=\small,text=AtlasInk] at (\v) {$\v$};
      \foreach \v in {b,c}
        \node[below right=2.5mm,font=\small,text=AtlasInk] at (\v) {$\v$};
      \node[above=2.5mm,font=\small,text=AtlasInk] at (f) {$f$};
      \foreach \v in {#3} {
        \fill[AtlasRubric!18!AtlasPaper] (\v) circle (4mm);
        \fill[AtlasInk] (\v) circle (1.8pt);
      }
      \ifstrequal{#1}{#2}{%
        \node[cop,minimum size=8mm,font=\scriptsize] at (#1) {C/R};
      }{%
        \node[cop] at (#1) {C};
        \node[robber] at (#2) {R};
      }
    \end{scope}
  \end{tikzpicture}%
}

% One declaration per position, reused by the chase and the proof.
\newcommand{\AtlasPursuitPosition}[1]{%
  \ifcase#1\relax
  \or\AtlasPursuitTree{a}{d}{b,c,d,e,f,g,h}{a/b,b/c,c/d,d/e,b/f,b/g,c/h}%
  \or\AtlasPursuitTree{b}{d}{c,d,e,h}{b/c,c/d,d/e,c/h}%
  \or\AtlasPursuitTree{b}{e}{c,d,e,h}{b/c,c/d,d/e,c/h}%
  \or\AtlasPursuitTree{c}{e}{d,e}{c/d,d/e}%
  \or\AtlasPursuitTree{d}{e}{e}{d/e}%
  \or\AtlasPursuitTree{e}{e}{}{}%
  \fi
}

\newcommand{\AtlasPursuitLegend}{%
  {\small C: cop\qquad R: robber\par
   Halos mark the robber's region.\par
   Dashed branches are inaccessible.\par}%
}

\newcommand{\AtlasPursuitBeat}[3]{%
  {\usebeamercolor[fg]{structure}\scshape\strut #1\strut\par}
  \medskip
  \parbox[t][2.3cm][t]{\linewidth}{\raggedright\strut #2\strut}\par
  \bigskip
  {\large Remaining vertices: $#3$\par}%
}

\begin{frame}{One cop is enough}
  % The local font change makes a statement slide without a new theme API.
  \begingroup
    \setbeamerfont{block title}{size=\large}
    \Large
    \begin{theorem}[Trees]\label{thm:tree}
      One cop can catch a robber\par
      on any finite tree.
    \end{theorem}
  \endgroup
  \bigskip
  The strategy: follow the unique path toward the robber.
\end{frame}

\begin{frame}{Follow the unique path}
  \framesubtitle{One legal chase on a fixed tree}
  \begin{columns}[T,onlytextwidth]
    \begin{column}{.54\textwidth}
      \begin{overlayarea}{\linewidth}{5.6cm}
        \centering
        \only<1|handout:0>{\AtlasPursuitPosition{1}}%
        \only<2|handout:0>{\AtlasPursuitPosition{2}}%
        \only<3|handout:0>{\AtlasPursuitPosition{3}}%
        \only<4|handout:1>{\AtlasPursuitPosition{4}}%
        \only<5|handout:0>{\AtlasPursuitPosition{5}}%
        \only<6|handout:0>{\AtlasPursuitPosition{6}}%
        \par\medskip
        \AtlasPursuitLegend
      \end{overlayarea}
    \end{column}
    \begin{column}{.43\textwidth}
      \begin{overlayarea}{\linewidth}{5.6cm}
        \raggedright
        \only<1|handout:0>{\AtlasPursuitBeat{1. Start}
          {The cop is at $a$; the robber is at $d$.
           Every other vertex is still reachable without crossing the cop.}{7}}
        \only<2|handout:0>{\AtlasPursuitBeat{2. The cop advances}
          {Move from $a$ to $b$. The branches through $a$, $f$, and $g$
           are now cut off from the robber.}{4}}
        \only<3|handout:0>{\AtlasPursuitBeat{3. The robber retreats}
          {The robber moves from $d$ to $e$. This changes the position,
           but cannot reopen a discarded branch.}{4}}
        \only<4|handout:0>{\AtlasPursuitBeat{4. Narrow the region}
          {The cop moves to $c$; the robber waits at $e$.
           The branch through $h$ is cut off too.}{2}}
        \only<5|handout:0>{\AtlasPursuitBeat{5. Close the last exit}
          {The cop moves to $d$. The robber stays at $e$:
           its only neighbor is occupied.}{1}}
        \only<6|handout:0>{\AtlasPursuitBeat{6. Capture}
          {The cop moves to $e$. Both pieces occupy the same vertex.
           There is nowhere left to retreat.}{0}}
        \only<beamer:0|handout:1>{%
          {\usebeamercolor[fg]{structure}\scshape The whole chase\par}
          \smallskip
          \small
          Step 4 is pictured.\par\smallskip
          \renewcommand{\arraystretch}{.9}
          \begin{tabular}{@{}crrc@{}}
            \toprule
            Step & C & R & Region \\
            \midrule
            1 & $a$ & $d$ & 7 \\
            2 & $b$ & $d$ & 4 \\
            3 & $b$ & $e$ & 4 \\
            4 & $c$ & $e$ & 2 \\
            5 & $d$ & $e$ & 1 \\
            6 & $e$ & $e$ & 0 \\
            \bottomrule
          \end{tabular}
          \smallskip\par
          R waits in steps 4 and 5.
        }
      \end{overlayarea}
    \end{column}
  \end{columns}
\end{frame}

\begin{frame}{Why every retreat ends}
  \framesubtitle{Return to the middle of the chase}
  \begin{columns}[T,onlytextwidth]
    \begin{column}{.54\textwidth}
      \begin{overlayarea}{\linewidth}{5.6cm}
        \centering
        \only<1|handout:0>{\AtlasPursuitPosition{3}}%
        \only<2-|handout:1>{\AtlasPursuitPosition{4}}%
        \par\smallskip
        \only<1|handout:0>{Before the move: $|S|=4$.}
        \only<2-|handout:1>{After the move: $|S'|=2$.}
        \par\smallskip
        {\small Cop: $b\to c$. Robber: stays at $e$.\par}
      \end{overlayarea}
    \end{column}
    \begin{column}{.43\textwidth}
      \begin{overlayarea}{\linewidth}{5.6cm}
        \small\raggedright
        {\usebeamercolor[fg]{structure}\scshape 1. Name the region\par}
        Delete the cop's vertex. Let $S$ be the robber's component.
        \par\smallskip
        \uncover<2->{%
          {\usebeamercolor[fg]{structure}\scshape 2. Move the boundary\par}
          Move to $v$ on the unique path. Every safe reply stays in
          \par\smallskip
          {\centering $\displaystyle S'\subseteq S\setminus\{v\}.$\par}
          \smallskip
        }
        \uncover<3->{%
          {\usebeamercolor[fg]{structure}\scshape 3. Count down\par}
          Since $v\in S$, $|S'|<|S|$.\par
          Finiteness forces capture.
          \hfill\usebeamertemplate{qed symbol}
        }
      \end{overlayarea}
    \end{column}
  \end{columns}
\end{frame}
